\documentclass[12pt]{report}

% FIX: "Option clash for package url" -- pdfx carica internamente hyperref, che a sua
% volta carica internamente url SENZA l'opzione hyphens. Se poi più giù si fa
% \usepackage[hyphens]{url}, url risulta già caricato con opzioni diverse -> clash.
% Soluzione: registrare l'opzione con \PassOptionsToPackage PRIMA di pdfx, cosi quando
% hyperref (dentro pdfx) carica url, lo fa gia con hyphens attivo. Non serve piu
% \usepackage[hyphens]{url} esplicito piu sotto.
\PassOptionsToPackage{hyphens}{url}
\usepackage{hyperref}
\usepackage[a-1b]{pdfx}
\usepackage{cjhebrew}
\usepackage{graphicx}
\graphicspath{{images/}}

\usepackage[a4paper, margin=2.5cm]{geometry} %Margini

\def\UrlBreaks{\do\/\do-\do_\do.} % Caratteri di spaziatura bibliografia

\usepackage[utf8]{inputenc}
\usepackage[T1]{fontenc}
\usepackage{csquotes}
\usepackage[british]{babel}

\usepackage{mathptmx}

\usepackage{algorithm} % Pseudocodice
\usepackage{algpseudocode} % Pseudocodice
% \floatname{algorithm}{Algoritmo} % Traduzione label

\usepackage[nottoc]{tocbibind} % Indice delle figure nel toc

\usepackage{multicol} % Due sezioni disposte per riga

\usepackage[backend=bibtex,sorting=none,block=space]{biblatex} % Bibliografia
\addbibresource{biblio.bib}

\usepackage{titlesec}
\usepackage{epigraph}
\epigraphnoindent
\setlength{\epigraphrule}{0.4pt}
\newcommand{\blankpage}{
    \newpage
    \thispagestyle{empty}
    \mbox{}
}

\titleformat{\chapter}
{\normalfont\huge\bfseries}
{\thechapter}
{1em}
{}

\titleformat{\section}
{\normalfont\LARGE\bfseries}
{\thesection}
{1em}
{}

\titleformat{\subsection}
{\normalfont\Large\bfseries}
{\thesubsection}
{1em}
{}

\usepackage{setspace}
\onehalfspacing
\usepackage{titlesec}
% \titleformat{\paragraph}
%   {\normalfont\bfseries}{}{0em}{}
\titlespacing*{\paragraph}
  {0pt}{1.2em}{0.6em}
  
% \setlength{\parindent}{0pt} % No tab a inizio paragrafo

\title{Bio-inspired learning for motion planning of collaborative robots}
\author{Davide D'Acunto}
\date{Ottobre 2026}

\usepackage{amsmath,amssymb}
\usepackage{booktabs}
\usepackage{enumitem}

\hypersetup{
	colorlinks=true,
	linkcolor=black,
	citecolor=blue,
	urlcolor=blue,
	pdftitle={Bio-inspired learning for motion planning of collaborative robots},
	pdfauthor={Davide D'Acunto}
}

\newenvironment{protocol}
{%
  \begin{minipage}{\textwidth}
  \begin{list}{}{%
    \setlength{\leftmargin}{0cm}%
    \setlength{\rightmargin}{0.8cm}%
    \setlength{\listparindent}{0pt}%
    \setlength{\itemindent}{0pt}%
  }
  \item\relax
  \begin{spacing}{0.9}
  \ttfamily\small
  \ignorespaces
}
{%
  \end{spacing}
  \end{list}
  \end{minipage}
}

\begin{document}

\numberwithin{figure}{chapter}
\numberwithin{table}{chapter}
\begin{titlepage}
    \centering
    % \vspace*{.5cm}

    \textbf{\Huge Università degli Studi di Salerno}

    \vspace{0.5cm}
    \huge Dipartimento di Ingegneria dell'Informazione ed Elettrica e Matematica Applicata
    
    \vspace{0.5cm}

    \LARGE Corso di Laurea Magistrale in Ingegneria Informatica (LM-32)

    \vspace{0.5cm}
    \includegraphics[width=170pt]{images/logo.png}

    \vspace{0.5cm}
    \huge Master Thesis in Natural Computation 

    \vspace{0.5cm}
    \textbf{\huge Bio Inspired Planning for Intentional Movements in Collaborative Robotics}

    \vspace{0.5cm}

    \begin{multicols}{2}
    \raggedright \textbf{\Large Supervisors:} \\
    \Large \raggedright Prof. Angelo Marcelli \\
    \Large \raggedright Prof. Pasquale Chiacchio \\

    \columnbreak
    \raggedleft \textbf{\Large Candidate:} \\
    \raggedleft \Large Davide D'\!Acunto \\
    \raggedleft \Large Num. 0622702524

    \end{multicols}
    \vspace{.5cm}
    Academic Year 2025/26

\end{titlepage}
\pagenumbering{gobble}
\null\vfill
\epigraph{Then what is your hand for?}{\textit{Rei Ayanami, End of Evangelion}}
\smallskip
\epigraph{To effectively contain a civilization's development and disarm it across such a long span of time, there is only one way: kill its science.
}{\textit{Liu Cixin, The Three-Body Problem}}
\smallskip
\epigraph{If one's different, one's bound to be lonely.}{\textit{Aldous Huxley, Brave New World}}
\smallskip
\epigraph{The saddest aspect of life right now is that science gathers knowledge faster than society gathers wisdom.}{\textit{Isaac Asimov, Book of Science and Nature Questions}}
\epigraph{\cjRL{'mt}}{\textit{The Golem of Prague}}
\smallskip
\epigraph{The true measure of a man is not his intelligence or how high he rises in this freak establishment. No, the true measure of a man is this: how quickly can he respond to the needs of others and how much of himself he can give.}{\textit{Philip K. Dick, Our Friends}}
\vfill
\blankpage
\chapter*{Acknowledgements}
\pagenumbering{Roman}
Write acknowledgements here.
\tableofcontents
\listoffigures
\listoftables

\clearpage

\pagenumbering{arabic}

% Content Outline

% 1. Introduction 3P
    % 1.1 Problem Definition
        % Description of collaborative robotics
        % Why the need for a human-like motion
        
    % 1.2 Problem Relevance in Computer Engineering
        % 

% 2 State of the Art 5-10P
    % 2.1 Current Research in Human-Robot Interaction
        % Studies about 

    % 2.2 Literature Limitations and Open Challenges
        % Knowledge of human motion in 3D space
        % Lack of purely 3D algorithms
        % 

% 3 Proposed Solution 20P
    % 3.1 Methodology
    % 3.2 System Design
        % 3.2.1 Hapti Device
    % 3.3 Haptic Interface Software
        % 3.3.1 ROS2 Controller
    % 3.4 Data Processing Pipeline
        % 3.4.1 Temporal Characterization
        % 3.4.2 Spatial Characterization
    % 3.5 Generative Model

% 4 Experimental Validation 20P
    % 4.1 Task Design
        % 4.1.1 Waypoint Definition
        % 4.1.2 Path Description
        % 4.1.3 Path Taxonomy
    % 4.2 Evaluation Metrics
    % 4.3 Experimental Protocol
        % 4.3.1 Source of Bias
        % 4.3.2 Tunable Parameters
        % 4.3.3 Parameter Values
        % 4.3.4 Written Protocol
    % 4.4 Experimental Data
    % 4.5 Statistical Significance
    % 4.6 System Evaluation

% 5 Conclusion 3-5P

% 6 Bibliography

%==============================================================================
% 1. Introduction
%==============================================================================
\chapter{Introduction}

Human motion, while natural and effortless for those who perform it, results from the coordination of the nervous, sensory and musculoskeletal systems, and has long been studied in biomechanics, medicine, neuroscience, robotics and computer science.

In robotics, it is particularly relevant to human--robot interaction, where collaborative robots share their workspace with humans. Humans are used to cooperating with other humans, and they read the movements of a partner through their own motor repertoire~\cite{bisio2014}. Robot motion that follows the kinematics of human movement has been shown to shorten the reaction times of human partners~\cite{huber2008}, to reduce the forces they apply during physical interaction~\cite{maurice2018} and to be perceived as more comfortable~\cite{coccaro}. A robot that moves like a human is therefore more predictable, and a predictable robot is easier to work with and more likely to be accepted, both on production lines and in everyday environments.

Predictability, however, is only part of the problem. Humans also reveal their intentions through their movements: the trajectory and velocity of a reaching hand indicate which object is about to be grasped. Making robot motion intent-expressive, or legible, is a related but distinct goal, since legible and predictable motion can contradict each other~\cite{dragan2013}. This work focuses on predictability, achieved through human-like motion, and leaves legibility to future developments.

Beyond human--robot interaction, models of human motion support the design of haptic interfaces, rehabilitation and assistive robots, and prosthetic devices, all of which benefit from a faithful description of how humans move.

This work has its roots in neuroscience, and adopts as its theoretical foundation the Kinematic Theory of rapid human movements~\cite{plamondon1995a,plamondon1995b}, in which a movement is described as the superimposition of elementary strokes, each with a lognormal velocity profile. The theory has been applied mostly to planar handwriting and signatures; showing that it can describe, and generate, movements in three-dimensional space extends its scope and may support future research on human motor control.

Human movements were recorded through an experimental campaign in which subjects moved the stylus of a haptic device through sequences of target points displayed on a screen, while the device attracted them towards each point without imposing the path. The velocity profile of each recording was decomposed into lognormal strokes, and the path of each stroke was described by a cubic Bézier segment, so that consecutive strokes form a continuous trajectory. Drawing the stroke parameters at random would not produce human-like movements, since the parameters are related to each other and to the task. A generative model was therefore trained on the extracted parameters to learn their distribution, conditioned on the target points, and to produce new movements from the same elementary units.

The contributions of this work are:
\begin{itemize}
  \item an experimental protocol and acquisition pipeline for recording three-dimensional movements through a haptic device;
  \item a representation of these movements as lognormal strokes in time and Bézier segments in space;
  \item a generative model that produces new human-like movements from this representation;
  \item a set of metrics to evaluate the human-likeness of the generated movements.
\end{itemize}

The remainder of this chapter defines the problem and discusses its relevance to computer engineering. Chapter~2 reviews the state of the art in human--robot interaction, neuroscience and the kinematics of human movement, which motivates the choices made in this work. Chapter~3 describes the proposed solution, its implementation and the reasons behind each choice. Chapter~4 presents the experimental validation, including the task design, the evaluation metrics, the experimental protocol and the results. Finally, Chapter~5 summarises the contributions and discusses future research.

\section{Problem Relevance in Computer Engineering}
As stated in the previous section, each task addressed in this project requires the application of knowledge and methodologies from the field of computer engineering. This involves identifying an appropriate theoretical model and translating it into a practical implementation through the application of established software design patterns. It also entails characterizing and understanding the specifications of a data acquisition device, as well as interfacing with its driver, given that the device itself is an electronic component.

Furthermore, the implementation of a generative neural network requires defining its architecture, selecting an appropriate loss function, and identifying the most effective representations of its inputs and outputs to optimize its performance. These activities collectively constitute fundamental aspects of artificial intelligence engineering.


%==============================================================================
% 2. State of the Art
%==============================================================================
\chapter{State of the Art}
The following section discusses the current advancements in human-robot interaction, and the limitations of the current literature. The goal is to provide the reader with a comprehensive overview of the state of the art in this field, highlighting challenges and opportunities for future research.

\section{Current Research}
As this work is traversal and requires fusing more areas of research together to reach the goal, the state of the art is then divided into main components, for the purpose of giving an overview in these fields.
\subsection{Human--Robot Interaction}
Collaborative robots~\cite{montini2024}, so-called cobots, interact with humans to execute a task in a shared workspace. Whether anthropomorphic visual cues aid cooperation is still debated~\cite{fink2012,onnasch2021,wittmann2024}, and appearance alone carries a risk: according to the uncanny valley hypothesis~\cite{mori2012}, a robot that closely resembles a human but is betrayed by some detail provides discomfort rather than affinity.

Motion, instead, has a measurable effect on interaction. When a robot handed objects to a human with a minimum-jerk velocity profile of the end-effector, rather than with a trapezoidal profile in joint space, the receivers reacted significantly faster~\cite{huber2008}. In such hand-overs, humans rely on both the position and the kinematics of the partner's movement to predict when and where the object will be delivered, so that a robot has to move predictably for joint action to succeed~\cite{glasauer2010}. Comfort can also be quantified: Coccaro et al. defined a Human-Likeness and Comfort index for robot trajectories, and found a consistent relation between perceived comfort and adherence to a lognormal velocity profile~\cite{coccaro}.

Human-like motion also engages the motor system of the observer. Observing an action automatically activates congruent motor representations, a mechanism known as motor resonance. A humanoid robot interferes with the movements of an observer when it moves with a minimum-jerk velocity profile, suggesting that a human-like velocity profile is sufficient for the robot to be perceived as an interaction partner~\cite{kupferberg2011}. Similarly, motor contagion, in which the observer unintentionally replicates features of the observed movement, appears with both human and robotic partners, unless the robot violates the biological laws of motion~\cite{bisio2014}. What matters to the observer is therefore not only how the robot looks, but whether its movement belongs to the observer's own motor repertoire.

These preferences persist when the interaction is physical. Humans following a robot along a curved path applied significantly lower forces when the robot moved with a biological velocity pattern, indicating that the features of natural movement remain a strong preference in guided movements~\cite{maurice2018}. Touch thus becomes a channel through which intentions are exchanged~\cite{hapticcomm}, and the predictability of the robot directly affects the effort required of its human partner.

Moving like a human, however, does not guarantee that the robot conveys its intentions. Dragan et al. distinguished predictability, the property of a motion that matches what an observer expects once the goal is known, from legibility, the property that allows the observer to infer an unknown goal from the motion, and showed that the two can contradict each other~\cite{dragan2013}. A human-like motion planner primarily addresses predictability, whereas legibility requires the robot to plan by explicitly modelling the inferences of the observer.

Human-like robot motion has been generated in several ways, from minimum-jerk interpolation in Cartesian space~\cite{huber2008} to dynamical movement primitives learned from demonstration~\cite{ijspeert2013}. This work follows a different route: robot motion is built from the lognormal strokes of the Kinematic Theory, described later in this chapter, so that the generated trajectories are composed of the same elementary units used to describe rapid human movements.

\subsection{Neuroscience}
Most knowledge on human kinematics and haptic communication stems from neuroscience. A voluntary movement originates from an intention rather than from an external stimulus~\cite[p.~816]{scott2021}. Its goal is probably set by the prefrontal cortex, the plan that realises it emerges from the interaction between the posterior parietal and the premotor cortices, and its execution is left to the primary motor cortex, the brainstem and the spinal cord~\cite[p.~710]{kandel2021}. Although this hierarchy is useful to describe the process, many of these stages operate in parallel~\cite[p.~710]{kandel2021}.

\paragraph{State estimation}
To plan a movement, the brain needs to know where the target is and where the limb is. Both estimates converge in the posterior parietal cortex, which identifies the targets of the action and estimates the state of the body~\cite[p.~857]{scott2021}. Visual information reaches the primary visual cortex V1 through the lateral geniculate nucleus of the thalamus. From V1, the dorsal stream, devoted to the visual guidance of movement, like reaching movements, arrives to both the middle temporal area (MT/V5), specialized in visual motion, and area V6, sensitive to the optic flow generated by self-motion~\cite{pitzalis2013}, before converging on the posterior parietal cortex~\cite[p.~505]{gilbertdas2021}.

Proprioception, the sense of the position and movement of the limbs, arises from muscle spindles, which signal muscle length and its rate of change, and from Golgi tendon organs, which signal muscle force~\cite[p.~421]{gardner2021}. These signals ascend through the spinal cord, cross to the opposite side in the brainstem and, after a relay in the thalamus, reach the primary somatosensory cortex, which passes them on to the posterior parietal cortex~\cite[pp.~450--452, 458]{gardner2021touch}. The perceived joint angle combines these afferents with the efferent motor command, so proprioception also reflects what the brain itself is commanding~\cite[p.~421]{gardner2021}. 

Vision and proprioception are fused, each weighted by its reliability: vision is more accurate in the horizontal direction, proprioception in depth~\cite[p.~720]{wolpert2021}. For this reason, the position of the limb can be estimated even when it is not in view. Since sensory signals arrive with a delay, the estimate is completed by a forward model that predicts the consequences of the current motor command from its efference copy~\cite[pp.~721--722]{wolpert2021}.

\paragraph{Motor plan}
The estimates of target and body state are passed to the premotor areas, which select the action and specify the spatiotemporal features of the movement~\cite[p.~710]{kandel2021}. The dorsal premotor cortex plans arm movements towards the target, the ventral premotor cortex shapes the hand for grasping, and the supplementary motor area is involved in self-initiated actions and in sequences of movements~\cite[pp.~829--830, 857]{scott2021}. The choice among competing actions is resolved by the basal ganglia, whose tonic inhibition is selectively released on the channel of the selected action~\cite[p.~940]{redgrave2021}.

Planning and execution are separable in time. When a movement is instructed but delayed, premotor neurons already encode it, and their activity during the delay predicts errors before the movement starts~\cite[p.~821]{scott2021}. During this preparation, the population activity evolves in a subspace that does not drive the muscles, and moves into one that does at movement onset~\cite[p.~834]{scott2021}.

The plan is expressed in task space rather than joint space: reaching movements follow nearly straight paths with smooth, bell-shaped velocity profiles that scale with distance, while the corresponding joint trajectories vary widely~\cite[pp.~725--726]{wolpert2021}. Similarly, a letter keeps its shape when written with different effectors, a property known as motor equivalence, which indicates that the plan is not tied to a specific set of muscles~\cite[pp.~726--727]{wolpert2021}. Among the infinite ways to reach a goal, the motor system appears to choose the one that minimizes a cost combining task success and effort~\cite[p.~727]{wolpert2021}.

\paragraph{Touch}
Interaction through a haptic device is perceived mainly through touch and proprioception. Four classes of mechanoreceptors in the hand decompose the contact with an object: rapidly adapting receptors signal contact, slip and vibration, while slowly adapting receptors signal pressure, grip force and skin stretch~\cite[pp.~437, 446--447]{gardner2021touch}. When a tool is used skilfully, the conditions at its working surface are perceived as if the fingertips were there, through the vibrations and forces transmitted by the tool~\cite[pp.~435--436]{gardner2021touch}. Touch is active when the subject moves the hand to explore, and passive when the stimulus is imposed from outside: the receptors are the same, but during active touch a copy of the motor command reaches the somatosensory pathways, allowing self-generated signals to be distinguished from external ones~\cite[p.~436]{gardner2021touch}.

This distinction has direct consequences for physical interaction. When a load is self-generated, grip force rises with it without delay, whereas an externally imposed load is followed with a lag and a higher grip force~\cite[pp.~723--724]{wolpert2021}. Similarly, a delay introduced by a robotic interface between a movement and its tactile consequence prevents the prediction from cancelling the self-generated sensation~\cite[p.~723]{wolpert2021}. When external forces cannot be predicted, the limb is stiffened by co-contracting antagonist muscles~\cite[pp.~716, 730]{wolpert2021}. Skilled movement nevertheless remains largely predictive, so that the loss of touch makes it clumsy rather than impossible~\cite[p.~464]{gardner2021touch}.

\paragraph{Online corrections}
Sensory feedback can correct a movement only after a delay: spinal reflexes respond within about 25~ms, transcortical responses within 50--100~ms and voluntary corrections within about 110~ms~\cites[p.~716]{wolpert2021}[p.~852]{scott2021}, while a complete visuomotor loop takes 120--150~ms~\cite[p.~717]{wolpert2021}. The early part of a rapid movement is therefore executed in open loop, and corrections can be introduced only once the feedback arrives. Feedback with high gain and long delay makes control unstable and oscillatory~\cite[p.~720]{wolpert2021}; two strategies limit this problem: intermittency, in which the movement is interrupted by pauses longer than the loop delay, as in manual tracking, and prediction~\cite[p.~723]{wolpert2021}. Moreover, the motor system does not correct every deviation, but only those relevant to the task, according to the minimal intervention principle~\cite[pp.~728--729]{wolpert2021}. Since motor noise grows with the size of the command, more accurate movements must be executed more slowly~\cite[pp.~727--728]{wolpert2021}. Finally, novel movements rely more on feedback than well-learned ones~\cite[pp.~729--732]{wolpert2021}, so corrective adjustments are expected to be more frequent early in learning. Variability is therefore not uniform: it is reduced along task-relevant dimensions and left free along the others~\cite[pp.~728--729]{wolpert2021}.

\paragraph{Motor learning}
How a movement is planned and executed depends on how well it is known. A novel movement relies on explicit strategies and on sensory feedback, whereas a well-learned movement is executed automatically through internal models, with feedback used mainly for adjustment~\cite[pp.~729--732]{wolpert2021}. Two subcortical structures drive this transition, each with a different learning signal. The basal ganglia learn by reinforcement: dopaminergic neurons signal whether the outcome of an action is better or worse than expected, modulating the probability and vigor of the actions that are selected~\cite[pp.~943--946]{redgrave2021}. With extended training, control shifts from goal-directed actions, sensitive to their outcome, to habits triggered by the context, which engage different territories of the striatum~\cite[p.~946]{redgrave2021}. The cerebellum instead learns from errors: climbing fibers signal the mismatch between the predicted and the actual movement, and the resulting synaptic changes gradually suppress the incorrect commands until the error disappears~\cite[pp.~924--925]{bastian2021}.

Building on these mechanisms and on motor equivalence, Senatore~\cite{senatore2012} hypothesized that a movement is stored both as an effector-independent sequence of target points and as an effector-dependent sequence of motor commands, and that skill learning proceeds in two stages. In the early stage, the basal ganglia acquire the order of the points to reach, while movements are slow, attention-demanding and guided by visual and proprioceptive feedback, as in a child copying letters with nearly straight point-to-point strokes. With practice, the cerebellum acquires the corresponding sequence of motor commands, and the whole sequence comes to be executed as a single, smooth and anticipatory behaviour~\cite[pp.~34--38, 56--58]{senatore2012}.

Learning where the target points lie involves several systems. Over a period of seconds, neurons in the dorsolateral prefrontal cortex maintain spatial representations, while the rehearsal of the location of objects involves the modulation of parietal, temporal and occipital areas by the frontal and premotor cortices~\cite[p.~1292]{shohamy2021}. During reaching, posterior parietal areas encode the target relative to the gaze and to the hand, and keep it active while the movement is delayed~\cite[pp.~825, 827]{scott2021}. When targets are presented on a display rather than in the workspace, an arbitrary mapping between visual and motor space has to be learned, an ability impaired by lesions of the dorsal premotor cortex~\cite[p.~829]{scott2021}. Furthermore, a flat display removes binocular disparity, one of the cues used to perceive depth~\cite[pp.~550--552]{gilbert2021}, and depth is already the direction in which vision is less reliable than proprioception~\cite[p.~720]{wolpert2021}.

Prior motor experience shifts this balance even for movements that have never been performed. A new behaviour is learned easily when it can be approximated by combining existing motor primitives~\cite[p.~734]{wolpert2021}, and practice with a variety of tasks sharing the same structure facilitates the learning of new tasks with that structure, while also improving feedback control~\cite{braun2009}. Experienced individuals therefore approach a new movement of a familiar class with more accurate predictions, as shown by expert cricket batsmen, who anticipate the bounce of the ball with shorter-latency saccades than amateurs~\cite[p.~723]{wolpert2021}. This advantage is, however, specific to the class of movements practised, since motor adaptation generalizes poorly to conditions that differ from the trained ones~\cite[p.~855]{scott2021}.

From a kinematic perspective, the quality of motor control has been related to how closely velocity profiles match ideal lognormals, a property that develops with learning and declines with ageing~\cite{plamondon2013}.
\subsection{Human trajectories}
Several studies have investigated human motion: Fitts' law~\cite{fitts1954} relates movement duration to the required accuracy, while the minimum-jerk model~\cite{flash1985} predicts the bell-shaped velocity profiles typical of reaching movements by assuming the nervous system minimizes jerk. Curvilinear movements instead obey the two-thirds power law~\cite{lacquaniti1983}, which links angular velocity to local curvature, while the isochrony principle states that movement duration is largely independent of trajectory extent~\cite{viviani1995}.

An alternative is offered by the Kinematic Theory of rapid human movements, in which a movement is described as the superimposition of elementary strokes each with a lognormal velocity profile, generated by the response of the neuromuscular system~\cite{plamondon1995a,plamondon1995b}; this is the idea behind the Sigma-Lognormal model~\cite{oreilly2009}. Working on this framework, Ferrer et al. proposed iDeLog~\cite{idelog}, an iterative dual spatial and kinematic parameter-extraction method that reconstructs 2D trajectories from their Sigma-Lognormal parameters, introducing the Signal-to-Noise Ratio, computed separately on the velocity profile and on the spatial path, as a metric for reconstruction quality. The framework was later extended to three-dimensional movements with iDeLog3D, which reconstructs 3D kinematics by linking virtual target points through temporally overlapping planar circumferences~\cite{idelog3d}.

More recently, this framework has been applied to physical human-robot interaction to define a Human-Likeness and Comfort index for robot trajectories: Coccaro et al. showed a consistent relation between perceived comfort and adherence to a lognormal velocity profile, validating the index against comfort judgments~\cite{coccaro}. 
%Tying these findings to the debate on anthropomorphic robot design~\cite{fink2012,onnasch2021,wittmann2024}, this thesis aims to develop an anthropomorphic motion planner to ease collaboration between humans and collaborative robots.

The Kinematic Theory also links these descriptions to the modular organisation of the motor system. In neuroscience, a motor primitive is a control module, such as the temporal profile of activation of a muscle or a synergy of muscles activated together, and the motor output is the weighted sum of the active primitives~\cite[p.~734]{wolpert2021}. Primitives have been sought both at the kinematic level, as strokes and submovements, and at the kinetic level, as force fields and muscle synergies, with a complex mapping between the two~\cite{flash2005}. At the muscular level, a fast single-joint movement shows a bell-shaped velocity profile produced by an agonist burst, followed by an antagonist burst and often by a second agonist burst~\cite[p.~755]{enoka2021}. The Kinematic Theory models this pair: each stroke is the output of a synergy of an agonist and an antagonist neuromuscular system, whose impulse responses converge to lognormals because each system consists of many coupled subsystems~\cite{plamondon2003}, and within this framework the stroke has been proposed as a movement primitive~\cite{woch2004}. Consistently, when the speed profiles produced by a neuromusculoskeletal model of the arm were reconstructed with the Sigma-Lognormal model, the number of lognormals matched the number of muscles receiving different excitation patterns, and an online correction appeared as an additional lognormal directed towards a new target point~\cite[pp.~55, 75]{parziale2016}. Nevertheless, the lognormal describes the output of the whole neuromuscular cascade rather than a neural module, so this correspondence remains a hypothesis at the kinematic level.

\section{Literature Limitations and Open Challenges}
The three lines of research reviewed in this chapter, namely human--robot interaction, neuroscience and the kinematics of human movement, converge on two hypotheses. First, a human partner interprets the movement of a robot through their own motor repertoire: observed actions activate congruent motor representations, and this resonance is preserved with a robot only when its motion follows biological velocity profiles~\cite{kupferberg2011,bisio2014}. A human-like motion would therefore be easier to understand and to anticipate. Second, the lognormal strokes of the Kinematic Theory may correspond to the motor primitives from which complex movements are built: each stroke is modelled as the output of an agonist--antagonist synergy~\cite{plamondon2003}, and the stroke has been proposed as a movement primitive~\cite{woch2004}.

Neither hypothesis has yet been examined in a unified way. The evidence on motor resonance comes mostly from visually mediated interaction and from minimum-jerk or similar velocity profiles~\cite{kupferberg2011,bisio2014}, whereas physical interaction, in which the human partner is guided by the robot, has received less attention~\cite{maurice2018}. The lognormal description has been used to assess the comfort of robot trajectories~\cite{coccaro}, while using it to generate them, rather than to evaluate them, is the focus of this work. The correspondence between lognormal strokes and muscle activity has been supported by simulation~\cite[pp.~55, 75]{parziale2016}, but not measured in humans, and most applications of the Kinematic Theory concern planar handwriting and signatures, its extension to three-dimensional movements being recent~\cite{idelog3d}. Finally, a human-like motion addresses the predictability of a robot rather than its legibility, which requires the observer to be modelled explicitly~\cite{dragan2013}.

This work addresses part of these gaps: it describes three-dimensional movements recorded through a haptic device as superimposed lognormal strokes, and uses the extracted parameters to train a generative model that produces new movements from the same elementary units. It does not test the neural or muscular basis of the strokes, nor their effect on a human partner in physical interaction. These questions could be addressed in future work by complementing the kinematic analysis with electromyography, to relate each stroke to the activity of agonist and antagonist muscles, with whole-arm motion capture, to relate the end-effector strokes to joint coordination, and with electroencephalography, which has already been used to relate the timing of a motor command to the Kinematic Theory~\cite{oreilly2013eeg}.

% Idelog che non cerca di rappresentare il comportamento motorio ma solo di minimizzare l'errroe di ricostruzione mediante più stroke possibili.

%==============================================================================
% 3. Proposed Solution
%==============================================================================
\chapter{Proposed Solution}

This chapter gives a comprehensive overview of the solution adopted to address the problem at hand, starting from the device chosen, going through the data collection and processing systems, then ending with the generative model and its performance evaluation. The goal is to provide the reader with a clear understanding of the choices, weighing pros and cons, so to justify the decisions made.
% The problem has been approached by combining software modules, each responsible for a specific task and independent with respect to the others. This solution allows to easily replace modules, speeding up the development process by allowing the testing of the single units of the system, and avoids coupling.

\section{Methodology}
As required by the problem to solve, the chosen methodological approach focuses on the data. This decision was natural, as having to put together various independent system means looking for a common interface for all of them. And in this case, that common interface is the data flow. 

This data is the very essence of the work, and its flow and change must be carefully observed. From the data acquisition, human motion characteristics must be present in its content. There may be errors from the collection device (measurement), errors from the collection itself (bias), errors from the software modules themselves (bugs), and finally, error from the generative model itself, due to a wrong architecture. 

Having in mind this flow, the methodology originates immediately, and asks for modularity. Someday, the data collection may be done through motion trackers, recording the motion of another part of the body. Meanwhile, instead of choosing one step adopted by this solution, someone else may try a different technique. Whatever is the reason behind the change, and the module changed, the flow of data and its meaning will always be the same.

Other than modularity, in this methodology the data flows through an ordered sequence of actions: a pipeline. While this aspect seems to regard more the system design, it is important to have in mind that also the decision of the collection method, the device it needs, the task assigned, and the experimental campaign themselves are steps that shape the data during its flow. And these decision may be changed later, but the pipeline must still work nonetheless.

\section{System Design}
The requirements that arose from choosing this methodology now apply to the composition of subsystem inside the overall system itself. Each subsystem corresponds to an interchangeable module; as required by the methodology adopted each is a step of the pipeline in charge of managing data. 

The three main subsystems are the following
\begin{itemize}
    \item The data collection system concerns recording the human motion data
    \begin{itemize}
        \item The experimental campaign, containing experimental protocol and designed tasks
        \item The data collection hardware and software
        \item The kind of data expected in input
        \item The collected dataset provided as output
        \item The interface with the user, which contains a set of tasks that must be formally defined to be also interpretable by a model
    \end{itemize}
    \item The data processing system is responsible for extracting the features that best describe the recorded data
    \begin{itemize}
        \item The interface with the data collection system, through accessing the collected dataset
        \item The codification of a theoretical model behind the reconstruction of the collected data through optimization problems
        \item The reconstructed dataset, which entries are the reformulation of the collected data through a set of feature, explained by the model adopted
    \end{itemize}
    \item The data generation system manages the generation of new human motion data
    \begin{itemize}
        \item The interface with the data processing system, through reading the reconstructed dataset
        \item The interface with the tasks, which interpretation must be the same of the human subject to not introduce bias in the dataset
        \item The neural network model architecture, its training and generation methods
        \item The training procedure defined by the loss function, the set of features to take as input, as recurrent input from the previous task, the expected outputs, and the data sets
    \end{itemize}
\end{itemize}

The following subsections provide a deeper insight into each system, without focusing on the implementation of them, which main modules are instead detailed in the next three sections, one for each subsystem. This more general overview serves to understand how the three system work in an architectural point of view.

Regarding the reasoning behind parameters or techniques adopted by these subsystems, their explanation is instead given in the next chapter.

\subsection{Data Collection System}
The data collection system is responsible for recording the human motion data. It is composed by the experiments themselves, the haptic device used to record the human motion, and the software that manages the communication between the experimenter, the haptic device and the human subject.
\paragraph{Haptic Device} 
The device adopted for the problem is \textit{Touch}, from \textit{3D Systems}. It is a haptic device with six degrees of freedom, allowing it to move in three-dimensional space and rotate around three axes. It is equipped with a stylus that can be held by the user, allowing them to interact with the device. The device is able to produce some linear forces, which are applied to the user's hand through the stylus. 

This kind of devices is commonly used in teleoperation, in which the user controls a remote robot by moving the stylus and feeling the forces applied by the robot. In this case the device is used only to collect data from the user's movements. 

\subsection{Data Processing System}
The data processing system is responsible for processing the collected data by extracting the features that describe it according to the theoretical model. This model is based on the theory of kinematics of rapid human movements, which describes human motion as a sequence of primitives, called lognormal strokes, that can be superimposed to generate complex movements. 

\subsection{Data Generation System}
The data generation system is responsible for generating new human motion data, based on the features extracted from the collected data, and the task to perform described in an experiment. 


\section{Haptic Interface Software}
The haptic interface software manages both the user interface of the haptic device and the communication with the robot.

From the side of the experimenter, the software component allows them to control the experiment execution, by selecting its configurations, like the path to show, the number of repetitions, the sampling frequency, and so on.

From the side of the human subject, the software component interfaces with them through a simulated environment that projects the position of the haptic device inside a virtual space. In this space, the path is represented as an ordered sequence of spheres, while the device is represented as a sphere that can be moved by the user through their interaction with the device.

\subsection{\texttt{ROS2} Controller}
The haptic device controller is in responsible for managing the communication between experiment, haptic device and subject through sending commands and receiving data, and of the experiment execution. 

The controller takes as input the configuration of the experiment, which includes the path to show, the number of repetitions, the size of the waypoints, the forces to apply and so on. From this information, during the execution it is able to communicate with the experiment through the keyboard, and with the haptic device through its command and state interfaces. The controller is also responsible for collecting the data from the haptic device, which includes the position, orientation, linear and angular velocities.

Furthermore, it is in charge of communicating with the subject through the screen via  the visualization of the path during the pre-training phase. In this case, the role of the controller is to communicate a series of colored waypoints, and the current position of the haptic device to the visualization software \texttt{RViz2}, which is responsible for rendering the scene on the screen. The controller applies the haptic feedback to the user, by computing the forces to apply to the haptic device based on the current position of the device and the position of the next waypoint, multiplied of a gain factor and the distance between the two.
% The software component aids the user in following the path by providing them with visual and haptic feedback. The former is given through the visualization of a colored sequence of spheres, in which the color indicates the progress along the path. The latter is given by applying attractive forces towards the next waypoint. Though, since these feedbacks would poison the data, they are disabled during the collection phase. This is because the goal of the experiment is to collect data that reflects the user's natural movements, without any external influence, while taking into account the user awareness of the path to follow. Nonetheless, the more accurate the user is in following the path, the more accurate the collected data will be, and thus the more accurate the generative model will be. Although, this goes in pair with the number of repetitions, paths and waypoints. This is a trade-off between the accuracy of the data, the naturalness of the user's movements and the possible sample poisoning by their fatigue.

\section{Data Processing Pipeline}
In the data processing pipeline the collected data is transformed into a representation that reflects the adoption of the theoretical model of human motion. This is done by finding a set of features that, when combined, can reconstruct the original data sufficiently well. In order to do so, a formal definition of the theoretical model is required, as its implementation is given in these software components.




% \subsection{Temporal Characterization}
% In particular, the crucial information that can be extracted from the dataset is the lognormal parameters of each rapid movement across two waypoints, which can be used to generate new movements that are similar to the original ones. 

% In literature there are already some software components that can extract the lognormal parameters from human movements, such as the \textit{idelog3d} software. Once this software is applied to the dataset, it will output a new dataset that contains the parameters of each rapid movement across two waypoints, which can be learned by the generative model. Once trained, the model can provide the velocity profile of a new movement across two waypoints through the application of the lognormal parameters.

% \subsection{Spatial Characterization}
% Regarding the spatial characterization of the movements, it would be handy to develop a software able to characterize the spatial properties of the movements, such as the starting and ending points, the curvature, the length and so on. In other words, the software should be able to extract the spatial properties of each rapid movement across two waypoints, which can be used to generate new movements that are similar to the original ones.

% A possible interpretation of the spatial properties of a movement is to consider it as a Bézier curve, actually a weighted spline Bézier curve, which is defined by a set of control points to which the curve is attracted. The control points are not necessarily the waypoints, while the weights are the parameters that define the attraction of the curve to the control points. The generative model could come with a set of control points(and eventually weights) that define the shape of the curve. In this case, the difference between waypoints and control points is that the former are the points that the model is supposed to reach, while the latter are the points to which the model's curve is attracted to.

% For instance, one possible spatial representation is the spline Bézier curve, which comes in handy since it doesn't always reach the target points, rather it tries to get as close as possible to them, which is what humans do when they move. The Bézier curve is defined by a set of control points, which can be used to define the shape of the curve. In order to teach a generative model to draw such lines, it is necessary to extract them from the actual path samples draw from the subjects movements.
% This software component should be able to provide the generative model with properties that are independent of the time characterization and explained by the scientific literature.

\section{Generative Model}
The generative model is trained with the dataset of movements, extended with the temporal and spatial characterizations, to learn the underlying principles of human motion. Once trained, the model can generate new movements that are similar to the original ones, by providing the path in the space and its time characterization through a velocity profile derived from the lognormal parameters. The generated movements can then be executed on the same system that collected the original data.

%==============================================================================
% 4. Experimental Validation
%==============================================================================
\chapter{Experimental Validation}
\section{Task Design}
The goal of the experiment campaign is relatively simple: the collection of a dataset of rapid and natural human movements. However, in order to obtain a dataset of sort, it is mandatory to provide the same task to each subject. Nonetheless, the same input must be grasped equally by different subjects. Thus, the deletion of ambiguity in task definition is imperative.

\subsection{Waypoint Definition}
As previously mentioned, a path is defined as an ordered sequence of points in the space, called waypoints, each defined by its position in the space. Visually, they are represented through spheres, which size is a trade-off between the accuracy of the collected data and the naturalness of the user's movements. However, even though the user should try to follow the path by crossing the waypoints, some deviation from them is to be expected, since no human movement ever perfectly follows a path. In the same way, a generative model should think of waypoints as a reference, akin to an attractor, rather than a strict constraint.

A path of two waypoints is simply a straight line between the two. On the other hand, a path of three waypoints is interpretable as a curve, since the user is supposed to follow the path by trying to reach the waypoints, but not strongly forcing themselves to cross and step on them: the result is a curve that starts from the first waypoint, gets near the second one, and ends almost on the third one. The same principle applies to paths of more than three waypoints, which are interpreted as a sequence of curves, which curvature is inversely related to the alignment of the waypoints.

The issue with such a representation is that it doesn't provide any information about the curvature of the path, neither about the depth of the waypoints in the space. 

\subsection{Path Description}

A path, in simple terms, is an ordered sequence of waypoints. This description would be sufficient for a model, since its perception of a three-dimensional space is mathematical and thus unmistakable. However, a human's perception of a three-dimensional space is not mathematical, and thus it is not as accurate as the model one. Therefore, it is necessary to provide the human subject with a more intuitive representation of the path.

A path could represent a common task that a human subject may perform in their daily life, such as moving an object from one place to another, or using a particular tool. The path should be designed to be as intuitive and natural as possible, eventually combining the geometrical description with a semantic one, ensuring that the subject grasps and performs it smoothly, avoiding unnatural movements that could poison the collected data.

\subparagraph{Visual \& Haptic Feedback} An intuitive representation should make use of the human natural instinct to merge different sensory inputs to form a more complete perception of the environment. In this case the two only classes of sensory inputs available are the visual and haptic ones, since they are the only two that are directly transformed into positional information in the perceived space by the motor cortex. 

The visual feedback is shown to the subject through the ordered sequence of spheres projected in a simulated three-dimensional environment in the screen. The depth perception is given by the perspective of the camera, which can be moved by the subject to change their point of view. 

The haptic feedback is given to the subject through the haptic device's actuators, which apply forces to the subject's hand to guide them towards the path. These forces are applied toward the next waypoint, and let the user stay on each waypoint to give them time to memorize the positions of the waypoints. The more the user's ability to follow the path increases, the less the haptic feedback is applied. More importantly, the rest time on the waypoints is zeroed, since the user must follow the path without stopping on the waypoints. This should happen some time before the haptic feedback is removed, in order to let the user get used to the fact that they are supposed to follow the path without stopping on the waypoints.

\subparagraph{Motor Representation} Both the haptic feedback and the visual feedback are not used during the recording phase, but rather in a previous phase called pre-training. During this step, the user is aided in following the path through the visual and haptic feedback, until a certain level of proficiency is reached. This allows the subject to obtain a representation of the path inside their motor cortex, which is a fundamental aspect of human motion. Once this representation is formed, the user starts to follow the path without any aid, in order to test this proficiency, the feedbacks are slowly removed, until they are completely disabled. The user is then asked to follow the path without any aid, while the system collects their movements. 

\subparagraph{Curvature Representation}
The only issue with representing a path only by a sequence of points is that it only takes into account direct links between waypoints. Humans rapid movements don't go straight from one point to the following one. Instead, they try to store momentum by not stopping on each point, but making curves during the execution. This aspect must be extracted from the dataset and learned by the generative model, since it is a fundamental aspect of human motion.

An idea to make this aspect stand out from the collected data could be the introduction of arrows in the visual environment, although this means crowding it with more information that could be distracting or hard to interpret for the user. Also, this requires defining other tunable parameters, since the arrows would have to be placed at a specific distance from the waypoint it starts from, with a specific length, at a specific size and orientation. However, by only specifying the waypoints, the user is free to interpret the path in their own way, and thus they can execute it in a more natural way, which is what the experiments aim to achieve. In the end, it is mandatory to explain to the user that they are supposed to follow the path by trying to reach the waypoints, but not strongly forcing themselves to cross and step on them, since this would lead to unnatural movements: the user should be encouraged to follow the path smoothly.

\subsection{Path Taxonomy}
The paths that will be shown to the subject during the experiment are grouped into the following categories: 
\begin{itemize}
    \item Burst: impulsive movements. Simple paths that help the subject getting used to the haptic device. 
    \begin{itemize}
        \item Straight line over the $XY, YZ$ or $XZ$ planes
        \item Semicircumference over the $XY, YZ$ or $XZ$ planes
    \end{itemize}
    \item Geometric: two- and three-dimensional round geometric shapes
    \begin{itemize}
        \item Helix, growing over the $X$ or $Y$ axes
        \item Spiral, in the $XY$ plane
        \item Ellipses, one in the $XY$ plane and the other in the $XZ$ plane
        \item Circle, one in the $XY$ plane and the other in the $XZ$ plane
    \end{itemize}
    \item Symbol: letters of an alphabet
    \begin{itemize}
        \item Cursive lowercase L $\left(\ell\right)$ in both the $XY$ and $XZ$ planes
        \item Print uppercase S in an oblique plane in both senses
        \item Cursive Greek Phi $\left(\varphi\right)$ in the $XY$ plane
        \item Print uppercase $N$ which segments differ in plane 
    \end{itemize}
\end{itemize}
\section{Evaluation Metrics}

\section{Experimental Protocol}
The real challenge is to design the experiments in such a way that there isn't difference between the information received by the subjects during the experiment and the one received by the generative model during its training phase.

\subsection{Sources of Bias}
To ensure consistency across different subjects, the experimental protocol should be designed to minimize any potential biases or confounding factors that could affect the results. For instance, it is necessary to simplify some complex terms by finding more intuitive and user-friendly ways to explain them, so to ensure user comprehension and engagement. 

Additionally, the protocol should be designed to minimize any potential sources of error or variability in the data collection process, such as differences in the user's physical abilities or environmental factors that could affect their movements. This last point is particularly important, since the experiment is designed to be executed in a laboratory environment, where the user is asked to sit on a chair and interact with the system through a haptic device. The laboratory environment should minimize any potential distractions or sources of noise that could affect the user's movements and the data collection process. 

\subparagraph{Posture} The environment must be designed to ensure the user's safety and comfort, by providing them with a comfortable chair and ensuring that the haptic device is properly calibrated and adjusted to their physical characteristics (i.e. dominant hand, arm length, height, etc.) to reduce any potential sources of discomfort or fatigue and increase comfort during the experiment. The more the user feels relaxed, the more natural their movements will be, and thus the more accurate the collected data will be. 

\subparagraph{Observer effect} The user is aware of the fact that they are being observed and that their movements are being recorded, which could affect their behaviour and introduce some bias in the collected data. To prevent this kind of bias, the user should be informed about the purpose of the experiment and the data collection process, and they should be given the opportunity to ask questions and provide feedback. This will help to ensure that the user feels comfortable and engaged in the experiment, and that they are motivated to perform their best. Of course, any reveal of the purpose of the experiment should be done in a way that doesn't influence the user's behaviour, and thus bias the collected data.

\subsection{Tunable Parameters}
The real challenge is to design the experiments in such a way that there isn't difference between the information received by the subjects during the experiment and the one received by the generative model during its training phase.  

% In this scenario, the users interact with a system that shows them an ordered sequence of waypoints in the space, namely a path, and they are asked to follow them with their hand. As they move, the system collects the trajectory of their movements, expressed as an ordered sequence of timestamped points with respect to a reference frame. For convenience, the reference frame is assumed to be the robot's base frame, and the points are expressed in Cartesian coordinates, each sampled at a fixed frequency.
\subparagraph{Parameters}
The session does have various parameters that need to be tuned in a trade-off between the accuracy of the collected data, their naturalness, the fatigue of the user, the length of the experiment and the amount of data collected. The parameters that can be tuned are:

\begin{itemize}

    \item $W_s$: The size of the spheres that represent the waypoints of the paths, which can be increased to make it easier for the user to follow the path, but at the cost of reducing the accuracy of the collected data. A bigger sphere would require less training and less precision from the user, but it would also make the collected data less accurate, since the user would be able to deviate more from the path without being penalized. On the other hand, a smaller sphere would require more training and more precision from the user, but it would also make the collected data more accurate, since the user would be forced to follow the path more closely. Although, this forceful guidance could lead to unnatural movements, unnatural to a human, which would be detrimental to the generative model. 
    \item $P_l$: The number of waypoints of a path, which can be increased to make the path more complex, at the cost of increasing the difficulty for the user to follow the path. A more complex path would require more training and heavier cognitive effort, while increasing the data collected. Conversely, a simpler path would require less training and less cognitive effort, while reducing the data collected. The complexity of the path needs to be balanced with the user's ability to follow it, to collect accurate and natural data.
    \item $P_r$: The number of repetitions for each path, which can be increased to collect more data, at the cost of increasing the fatigue of the user and the length of the experiment. This parameter comes at the cost of making brief rest periods between repetitions, which would compensate the user's fatigue, but at the cost of greatly increasing the time duration of the full experiment session.
    \item $T_r$: The number of execution prior to the recording phase, named "pre-training phase". Albeit it would be simpler to have a fixed number of repetitions, the issue is within the variability of human learning. Some subjects may learn a path in less time than other, thanks to a prior knowledge of environments over the screen, or of their . Since this parameter depends on the user's ability to follow the path, and on the complexity of the path itself, it is left to decide up to the experimenter. If the user correctly follows the path for a certain number of repetitions, and feels enough confident to proceed to the recording phase, they can switch to the recording phase. Until then, the pre-training will aid the participant by providing the haptic feedback.
    

    \item $N_P$: The number of paths, which can be increased to collect more diverse data, at the cost of increasing the length of the experiment. More paths would provide a better representation of the task space and allow to more generalization, but would also require more time and effort from the user. Although, having a few paths would limit the diversity of the data collected, and thus the generalization of the generative model.
    \item $R_t$ and $R_f$: The rest periods length and frequency (expressed as a fraction of $N_P$), which can be increased to reduce the fatigue of the user, at the cost of increasing the duration of the experiment. The rest periods need to be short to keep the user focused on the task, but long enough to allow them to recover from their fatigue, while not being felt coercive by the user. Moreover, the rest periods need to be frequent enough to allow the user to recover from their fatigue while executing the different paths, but not too frequent to make the user lose focus on the task. 
\end{itemize}

\subsection{Parameter Values}
Regarding the aforementioned experimental parameters, the following values have been chosen: 

\begin{table}[h!]
    \centering
    \begin{tabular}{|ccc|}
    \hline
    \textbf{Parameter} & \textbf{Description} & \textbf{Value} \\
    \hline
    \hline
    $W_s$ & Waypoint diameter & 0.05 m \\
    $P_l$ & Path length & 2-7 waypoints \\
    $P_r$ & Path repetitions & 20 repetitions \\
    $N_P$ & Number of paths & 20 paths \\
    $R_t$ & Rest time & 60s \\
    $R_f$ & Rest frequency & $1/6$ \\
    \hline
\end{tabular}
\end{table}
As explained before, the training repetitions are not a tunable parameter, since they depend on the subject ability to follow the path, and on the complexity of the path itself. At most, it would be possible to define an acceptance criterion during the training phase. Since the training requires them to cross all the waypoints, it is safe to assume that the criterion is in some way already encoded in the waypoint size $W_s$.

\subsection{Written Protocol}
Providing the user with a written protocol is mandatory not only to ensure that all the user have the same information, but also to ensure that they have the same understanding of the task. From the different equivalent sequences of words available to explain an instruction, the most intuitive should be chosen, to avoid any bias in the collected data. 
\subparagraph{Input} In each experiment of a session a path is shown to the user, their goal is to traverse close to the waypoint as fast as possible, with the understanding that it's not strictly necessary to cross all the waypoints, but rather to think of them as attractors. They will first be pre-trained to follow the path with the aid of visual and haptic feedback, until they reach a certain level of proficiency, defined by the experimenter, and then they will be asked to follow the path without any aid, while the system collects their movements.
\subparagraph{Result} For each execution of a path, the system collects a sequence of timestamped points in the space, which are expressed in Cartesian coordinates with respect to the reference frame. Furthermore, the tangential velocity of the user's hand is computed at each timestamp. By interpolating these points, and by having a high sampling frequency, it is possible to reconstruct the trajectory of the user's movements with better accuracy. Other information like the user's hand orientation and the angular twist velocity are collected as well, but they are not needed for the generative model, since they are not used to generate new movements. They are collected to provide a more complete representation of the user's movements, and to eventually allow for a more detailed analysis of the data.

\begin{protocol}
... Welcome ... What's your dominant hand? (If no preference is expressed, pick left hand when possible). (Prepare the environment on the subject choice).

The device to your (left/right) side is a haptic device, which is connected to the screen in front of you (the screen shows an empty environment). In the screen, you can observe the effects that manipulating the device has on the environment.

Please extract the device from its housing and handle it as if it was a pen or a precision tool. Feel free to explore the environment in the simulation. Try to exit from it: you will feel forces that will prevent you from leaving it. Those are the limits within you can freely move. There are two buttons on the pen: the darker one lets you rotate the environment, the lighter one to zoom. [mouse]

(Once the subject learnt how to use it). Now, the screen will show you some colored spheres (inside other transparent spheres?). Try to make a smooth, fluid and natural movement that crosses all the spheres. The order in which the movement must cross the spheres is given by their number. 

Initially, when you try to cross the spheres, the device will help you reach them by pulling you towards their center. After a while, the device will let you traverse to the next sphere, with a slight pulling force.

After a fixed amount of repetitions, the screen won’t show any information, and only the haptic feedback will be left, which pulling force will decrease execution by execution.

After a fixed amount of repetitions, the pulling force will decrease to zero (it must dynamically adjust itself by studying the subject accuracy), from there you will be in the recording phase.

(Should the user be told when they're in the recording phase?)

Every time you want to be placed to the start position and/or reset the execution, just press the dark button. Now please, try to cross the spheres with the aid of the device.

[steps: (sight + haptic feedback), (haptic feedback), (no info)]

(When the pre training phase ends) Congrats. Now, proceed to repeat the same movement you've done before, without the screen information (or with only the points visible?) as accurately as possible. If you feel you're off track, don't fix the movement and try again, since it would be discarded anyways. (In order to help you, the device will automatically place itself in the starting position when you press the lighter button. A brief vibration will tell you when to start the movement. If during the movement you feel the same vibration, then you have to redo it again?)
\end{protocol}


Those are the same inputs and outputs of the generative model, with the exception that it won't follow the same policies, but only understand the underlying principles from the dataset.


% In the end, it is safe to assume that the parameters that can be tuned are $W_s$, $P_l$, $P_r$, $N_P$, $R_t$ and $R_f$. Another parameter would've been the sampling frequency, but it was set to the maximum possible frequency allowed by the hardware in order to collect the higher amount of data possible. In the end, the amount of short movements collected are the following: \begin{equation}
%     N_M = N_P \cdot P_r \cdot \overline{P_l}, \qquad \overline{P_l} = \frac{1}{N_P} \sum_{i=1}^{N_P} P_{l_i}
% \end{equation}
% On the other hand, the time length of the experiment is given by the following equation:
% \begin{equation}
%     T_E = P_r \cdot T_{P_i} + R_f \cdot R_t + T_0, \qquad \overline{T_{P_l}} = \frac{1}{N_P} \sum_{i=1}^{N_P} T_{P_{l_i}}
% \end{equation}
% Where $T_0$ is the time required to set up the experiment and introduce it to the human subject, and $T_{P_i}$ is the time required to execute the $i$-th path, taking into account the training time, which is given by the following equation:
% \begin{equation}
%     T_{P_i} = P_{l_i} \cdot W_s \cdot L_{i}
% \end{equation}
% Where $L_i$ refers to the workload of the user during the execution of the $i$-th path. It is not a value that can be directly extracted through a math equation, though it is possible to split it into components: the overall effort used in the experiment session, the effect of the overall time and the breaks on the user. The overall effort is given by the following equation:
% \begin{equation}
%     L_{i} = 
% \end{equation}

\section{Experimental Data}
Through the executions of the experiments, the system collects a dataset of human movements, expressed as a sequence of timestamped points in the space, with respect to the reference frame. 

\section{Statistical Significance}

\section{System Evaluation}

%==============================================================================
% 5. Conclusion
%==============================================================================
\chapter{Conclusion}
Write the conclusion here.

%==============================================================================
% 6. Bibliography
%==============================================================================
\printbibliography[heading=bibnumbered]
\end{document}
